% 卷、部分与目录共享全书计数；新卷不重置部分号和章号。
% 目录卷标题预留与下属条目同页的空间，避免空部分组跨页散开。
\usepackage{needspace}
\newcounter{volume}
\renewcommand{\thevolume}{\arabic{volume}}
\newcommand{\BookVolumeName}{第\chinese{volume}卷}
\newcommand{\BookVolumeTitle}{}
\newcommand{\BookPartEnglishTitle}{}
\ifBookVolumeEdition
  \newcommand{\BookTopBookmarkLevel}{-1}
\else
  \newcommand{\BookTopBookmarkLevel}{-2}
\fi
\newcommand{\BookTopBookmark}[2]{%
  \pdfbookmark[\BookTopBookmarkLevel]{#1}{#2}}

\NewDocumentCommand{\BookVolumeRef}{m m}{%
  \ifBookVolumeEdition #2\else 第~\ref{#1}卷\fi}
\NewDocumentCommand{\BookPartRef}{m m}{%
  \ifBookVolumeEdition #2\else 第~\ref{#1}部分\fi}
\NewDocumentCommand{\BookChapterRef}{m m}{%
  \ifBookVolumeEdition #2\else 第~\ref{#1}章\fi}
\NewDocumentCommand{\BookSectionRef}{m m}{%
  \ifBookVolumeEdition #2\else 第~\ref{#1}节\fi}
\NewDocumentCommand{\BookTheoremRef}{m m}{%
  \ifBookVolumeEdition #2\else 定理~\ref{#1}\fi}
\NewDocumentCommand{\BookEquationRef}{m m}{%
  \ifBookVolumeEdition #2\else 式~\eqref{#1}\fi}

\makeatletter
% 注册高于 part 的目录及 PDF 书签层级，不改动第三方模板。
\def\ttll@volume{-2}
\def\toclevel@volume{-2}
\def\ttll@bookbackmatter{-1}
\def\toclevel@bookbackmatter{-1}
\providecommand*{\theHvolume}{\arabic{volume}}
\def\Hy@numberline#1{%
  \ifnum\Hy@level=-2
    第\zhnumber{#1}卷\space
  \else\ifnum\Hy@level=-1
    第\zhnumber{#1}部分\space
  \else\ifnum\Hy@level=0
    第 #1 章\space
  \else
    #1\space
  \fi\fi\fi}
\newcommand{\BookBackmatterTocEntry}[1]{%
  \begingroup
    \long\def\Hy@writebookmark##1##2##3##4##5{}%
    \addcontentsline{toc}{bookbackmatter}{\protect\booktocbackmatterbadge{#1}}%
  \endgroup}
\newcommand{\BookBackmatterChapter}[2]{%
  \chapter*{#1}%
  \BookTopBookmark{#1}{#2}%
  \markboth{#1}{}%
  \BookBackmatterTocEntry{#1}}
\newcommand{\BookBackmatterIndexChapter}[1]{%
  \BookOriginalChapter*{#1}%
  \BookTopBookmark{#1}{book-index}%
  \markboth{#1}{}%
  \BookBackmatterTocEntry{#1}}
\renewcommand{\tableofcontents}{%
  \chapter*{\contentsname}\markboth{\contentsname}{}%
  \BookTopBookmark{\contentsname}{book-contents}%
  \begin{fullwidth}
    \@starttoc{toc}%
  \end{fullwidth}}
\makeatother

\contentsmargin{0pt}
\titlecontents{volume}[0pt]
  {\Needspace{60mm}\addvspace{17pt}\contentsmargin{0pt}%
    \heiti\large\mdseries\color{white}\hypersetup{linkcolor=white}}
  {\booktocvolumeband{\thecontentslabel}}{}%
  {\hfill\booktocpage}
  [\vspace{5pt}]
\titlecontents{part}[\booktoctitleindent]
  {\addvspace{9pt}\heiti\bfseries\color{BookTeal}}
  {\contentslabel[\booktocpartbadge{\thecontentslabel}]{25mm}}{}%
  {\hspace{0.6em}\titlerule*[0.6pc]{\textcolor{BookRule}{.}}\booktocpage}
\titlecontents{bookbackmatter}[\booktoctitleindent]
  {\addvspace{9pt}\heiti\bfseries\color{BookTeal}}
  {}{\hspace*{-25mm}}%
  {\hspace{0.6em}\titlerule*[0.6pc]{\textcolor{BookRule}{.}}\booktocpage}
\titlecontents{chapter}[\booktoctitleindent]
  {\addvspace{5pt}\heiti\color{BookInk}}
  {\contentslabel[第\thecontentslabel 章]{22mm}}{}%
  {\hfill\booktocpage}
\titlecontents{section}[\booktoctitleindent]
  {\small\color{BookMuted}}
  {\contentslabel{22mm}}{}{\hfill\booktocpage}

% 卷首页只接收中英文标题；具体视觉样式留在 layout.tex 中。
\makeatletter
\newcommand{\bookvolume}[2]{%
  \ifBookVolumeEdition\else
    \clearpage\thispagestyle{empty}%
  \fi
  \refstepcounter{volume}%
  \renewcommand{\BookVolumeTitle}{#1}%
  \NR@gettitle{#1}%
  \markboth{\BookVolumeName\quad #1}{}%
  \ifBookVolumeEdition\else
    \addcontentsline{toc}{volume}{\protect\numberline{\thevolume}#1}%
    \null
    \bookvolumetitlepage{#1}{#2}\par
  \fi
  \ignorespaces}

\newcommand{\bookpart}[2]{%
  \renewcommand{\BookPartEnglishTitle}{#2}\part{#1}}
\RenewDocumentCommand{\part}{s o m}{%
  % 不依赖 openright 或单双面选项：部首页始终使用奇数书页码。
  % 在建立标签、目录与书签锚点之前补空白页，避免链接落到空白页。
  \clearpage
  \ifodd\value{page}\else
    \null\thispagestyle{empty}\newpage
  \fi
  \thispagestyle{empty}
  \IfBooleanTF{#1}{\phantomsection}{\refstepcounter{part}}
  \NR@gettitle{#3}%
  \IfBooleanF{#1}{\IfNoValueTF{#2}
    {\addcontentsline{toc}{part}{\protect\numberline{\thepart}#3}}
    {\addcontentsline{toc}{part}{\protect\numberline{\thepart}#2}}}
  \markboth{#3}{}\null
  \begin{tikzpicture}[remember picture,overlay]
    \fill[BookPaper] (current page.north west) rectangle (current page.south east);
    \fill[BookTeal] ([xshift=23mm,yshift=-54mm]current page.north west) rectangle ++(57mm,-22mm);
    \node[anchor=center,text=white,font=\heiti\fontsize{21}{27}\selectfont]
      at ([xshift=51.5mm,yshift=-65mm]current page.north west)
      {\IfBooleanF{#1}{\BookPartName}};
    \node[anchor=north west,inner sep=0pt,text=BookInk,text width=169mm,
      font=\heiti\bfseries\fontsize{29}{38}\selectfont]
      at ([xshift=23mm,yshift=-106mm]current page.north west) {#3};
    \node[anchor=north west,inner sep=0pt,text=BookMuted,
      font=\sffamily\fontsize{11}{16}\selectfont]
      at ([xshift=23mm,yshift=-126mm]current page.north west) {\BookPartEnglishTitle};
  % 到下一个结构命令才输出，使随后的 label 指向当前入口页。
  \end{tikzpicture}\par\ignorespaces}
\makeatother
